\documentclass[11pt]{article}
\usepackage[margin=1in]{geometry}
\usepackage{amsmath,amssymb,amsthm}
\usepackage{xurl}
\usepackage{hyperref}
\sloppy
\let\oldus\_ \renewcommand{\_}{\oldus\allowbreak}
\newtheorem{theorem}{Theorem}
\newtheorem{lemma}[theorem]{Lemma}
\theoremstyle{definition}
\newtheorem{definition}[theorem]{Definition}
\title{Conjecture 00000003905 is false:\\ $\mu\bigl((x)^{[p^e]}\bigr)=1$ in $k[x,y]$, while $\dim R-\operatorname{ht}(x)=1$}
\author{Jackmeson1}
\date{2026-10-05}
\begin{document}
\maketitle

\section{The conjecture}
Conjecture 00000003905 (English text, verbatim): ``Definition: The Frobenius power $I^{[p^e]}$ is
the ideal generated by the $p^e$-th powers of the generators of $I$, and mu denotes the minimal number
of generators. Conjecture: For fixed $I$, the number $\mu(I^{[p^e]})$ is a polynomial in $p^e$ once
$e \ge e_0$, of degree equal to the ring dimension minus the height of $I$, and its leading
coefficient depends only on the initial ideal.'' The Chinese text says the same.

\paragraph{Reading.} $R$ is a Noetherian ring of prime characteristic $p$; since the text speaks of an
initial ideal, $R$ is a polynomial ring over a field. ``A polynomial in $p^e$ once $e\ge e_0$ of degree
$\dim R-\operatorname{ht} I$'' means: there are $e_0$ and a polynomial $P$ with
$\mu(I^{[p^e]})=P(p^e)$ for all $e\ge e_0$ and $\deg P=\dim R-\operatorname{ht}I$. We take the
coefficients of $P$ in an arbitrary field $K$ of characteristic $0$ (for instance $\mathbb Q$ or
$\mathbb R$), which covers integer and rational coefficients. We refute the degree clause, and hence
the conjunction, for $R=k[x,y]$ ($k$ any field of characteristic $p$, any prime $p$) and $I=(x)$.

\emph{Not refuted:} the bare polynomiality clause (it holds here: the value is constantly $1$); the
leading-coefficient clause (not addressed); readings restricted to ideals with
$\operatorname{ht} I=\dim R$ (for example $\mathfrak m$-primary ideals, where the predicted degree is
$0$); the local ring $k[x,y]_{(x,y)}$ (the same computation applies, but it is not formalized).

\section{Definitions}
\begin{definition}
For an ideal $I$ of a commutative ring $R$ and $q\in\mathbb N$, $I^{[q]}$ is the ideal generated by
$\{a^q : a\in I\}$ (Lean: \texttt{frobeniusPow}). Lemma~\ref{lem:gens} shows that for $q=p^e$ in
characteristic $p$ it equals the ideal generated by the $q$-th powers of the elements of \emph{any}
generating set of $I$, so it agrees with the definition in the conjecture and does not depend on the
chosen generators.
\end{definition}
\begin{definition}
$\mu(J)$ is the least cardinality of a generating set of $J$ (Mathlib: \texttt{Submodule.spanFinrank},
defined as the least cardinality of a generating set; it is set to $0$ only when no finite generating
set exists, which cannot happen in a Noetherian ring).
$\dim R$ is the Krull dimension (Mathlib \texttt{ringKrullDim}, the supremum of lengths of chains of
prime ideals) and $\operatorname{ht} I$ is the height (Mathlib \texttt{Ideal.height}, the infimum of the
heights of the minimal primes over $I$).
\end{definition}

\section{Proof}
\begin{lemma}[\texttt{frobeniusPow\_eq\_span\_pow\_gens}]\label{lem:gens}
Let $R$ have characteristic $p$, $q=p^e$, $S\subseteq R$. Then $(S)^{[q]}=(s^q : s\in S)$.
\end{lemma}
\begin{proof}
$\supseteq$ is clear. For $\subseteq$ it suffices that $a^q\in(s^q:s\in S)$ for every $a\in(S)$; by
induction on $a$: $0^q=0$; $s^q$ is a generator; $(a+b)^q=a^q+b^q$ in characteristic $p$; and
$(ra)^q=r^qa^q$.
\end{proof}

\begin{lemma}[\texttt{ringKrullDim\_eq\_two}, \texttt{height\_I}]
$\dim k[x,y]=2$ and $\operatorname{ht}(x)=1$.
\end{lemma}
\begin{proof}
Mathlib's \texttt{MvPolynomial.ringKrullDim\_of\_isNoetherianRing} states
$\dim R[X_1,\dots,X_n]=\dim R+n$ for Noetherian $R$; with $\dim k=0$ this gives $2$.
Mathlib's \texttt{Ideal.height\_span\_singleton\_eq\_one\_of\_mem\_nonZeroDivisors} (docstring: ``In a
Noetherian ring, the height of a principal ideal spanned by a non-unit non-zero-divisor is one'')
applies to $x$: it is nonzero in the domain $k[x,y]$, and it is not a unit because evaluation at $0$
sends it to $0$.
\end{proof}

\begin{lemma}[\texttt{mu\_frobeniusPow\_I}]
$\mu\bigl((x)^{[p^e]}\bigr)=1$ for every $e\ge 0$.
\end{lemma}
\begin{proof}
By Lemma~\ref{lem:gens}, $(x)^{[p^e]}=(x^{p^e})$. A nonzero principal ideal has $\mu=1$
(Mathlib \texttt{Submodule.spanFinrank\_singleton}), and $x^{p^e}\ne 0$.
\end{proof}

\begin{lemma}[\texttt{eq\_one\_of\_eval\_pow}]
If $P\in K[t]$, $p\ge 2$ and $P(p^e)=1$ for all $e\ge e_0$, then $P=1$.
\end{lemma}
\begin{proof}
The points $p^e$ ($e\ge e_0$) are pairwise distinct in $K$ (characteristic $0$), so $P-1$ has
infinitely many roots.
\end{proof}

\begin{theorem}[\texttt{conjecture\_3905\_false}]
For every prime $p$, field $k$ of characteristic $p$ and field $K$ of characteristic $0$: in
$R=k[x,y]$ with $I=(x)$, $\dim R=2$, $\operatorname{ht}I=1$, and there are no $e_0$ and $P\in K[t]$ with
$\mu(I^{[p^e]})=P(p^e)$ for all $e\ge e_0$ and $\deg P+\operatorname{ht}I=\dim R$.
\end{theorem}
\begin{proof}
By the lemmas, any such $P$ equals $1$, so $\deg P+\operatorname{ht}I=0+1=1\ne2$.
\end{proof}
We state the degree condition as $\deg P+\operatorname{ht}I=\dim R$ (in Mathlib's
$\mathrm{WithBot}\,\mathbb N_\infty$), which is equivalent to $\deg P=\dim R-\operatorname{ht}I$
because all three quantities are finite here. Here $\deg$ is \texttt{Polynomial.natDegree};
\texttt{frobenius\_mu\_polynomial\_is\_one} also proves $P=1$ and \texttt{Polynomial.degree} $P=0$, so the
conclusion does not depend on the degree convention.

\section{Lean formalization}
File \texttt{lean/Conjecture3905/Basic.lean} (151 lines, \texttt{import Mathlib} only; Lean
\texttt{v4.33.1}, Mathlib \texttt{v4.33.1}). Main theorems: \texttt{C3905.conjecture\_3905\_false},
\texttt{C3905.frobenius\_mu\_polynomial\_is\_one}; supporting:
\texttt{frobeniusPow\_eq\_span\_pow\_gens}, \texttt{ringKrullDim\_eq\_two}, \texttt{height\_I},
\texttt{frobeniusPow\_I}, \texttt{mu\_frobeniusPow\_I}, \texttt{eq\_one\_of\_eval\_pow}.
\texttt{lake build} succeeds; \texttt{\#print axioms} for both main theorems reports exactly
\texttt{[propext, Classical.choice, Quot.sound]}. There is no \texttt{sorry}, no \texttt{native\_decide}
and no added axiom.

\section{Relation to the earlier submission}
PR \#286 (by orionsheep) used $I=(x,y)$ in $k[x,y,z]$, with $\mu=2$ for every $q$. It was not merged:
its main theorem was a conjunction of numeral identities such as $(2:\mathbb N)=2$ and $2\cdot2-2=2$, and
no ring, ideal, Frobenius power or $\mu$ was defined in Lean. The reviewer asked for a version that
``defines $\mu$ on ideals in Lean''. This submission defines the Frobenius power in Lean and proves that
it does not depend on the generators. It uses Mathlib's $\mu$, Krull dimension and height, proves
$\dim k[x,y]=2$, $\operatorname{ht}(x)=1$ and $\mu((x)^{[p^e]})=1$ for every $e$, and proves that no
interpolating polynomial has degree $\dim R-\operatorname{ht}I$. We switched to the principal ideal $(x)$
because $\mu=1$ is immediate there, whereas $\mu((x^q,y^q))=2$ would need a separate non-principality
proof.

\section*{Sources}
\begin{itemize}
\item Mathlib, files \nolinkurl{RingTheory/KrullDimension/Polynomial.lean},
\nolinkurl{RingTheory/Ideal/KrullsHeightTheorem.lean} and \nolinkurl{Algebra/Module/SpanRank.lean}.
\item PR \#286: \url{https://github.com/The-Last-Math-Competition/The-Last-Math-Competition/pull/286}.
\end{itemize}
\end{document}
